\documentclass[10pt,a4paper]{amsart}
\usepackage[margin=19mm]{geometry}
\usepackage{amsmath,amssymb,mathtools}
\usepackage[scr=rsfs,cal=euler]{mathalfa}
\usepackage{xfrac}
\usepackage[unicode,hidelinks,bookmarks=false]{hyperref}
\newcommand{\ie}{i.e.\ }
\newcommand{\cf}{cf.\ }
\newcommand{\resp}{respectively\ }
\newcommand{\Subsection}{\subsection}
\providecommand{\nobreakdash}{\nobreak\hbox{-}}
\setlength{\emergencystretch}{2em}
\multlinegap0pt
\usepackage{bm}
\usepackage{bbm}
\usepackage{dsfont}
\usepackage{xspace}
\usepackage{cancel}
\usepackage{mathtools}
\usepackage{amsthm}
\makeatletter
\@ifundefined{proposition}{\newtheorem{proposition}{Proposition}}{}
\@ifundefined{lemma}{\newtheorem{lemma}[proposition]{Lemma}}{}
\makeatother
\DeclareMathOperator{\Frac}{Frac}
\newcommand\bU{{\mathbf U}}
\newcommand\cO{{\mathcal{O}}}
\newcommand\fp{{\mathfrak{p}}}
\title[Unramified Grothendieck-Serre for simply-connected group sch]{Unramified Grothendieck-Serre for simply-connected group schemes satisfying an isotropy condition via unipotent chains}
\author{Roman Fedorov}
\date{Source: arXiv:2204.05442v7 (CC BY 4.0)}
\begin{document}
\maketitle

\noindent\emph{Source and licence.} Unramified Grothendieck-Serre for simply-connected group schemes satisfying an isotropy condition via unipotent chains. Version arXiv:2204.05442v7 (CC BY 4.0), distributed under the Creative Commons Attribution 4.0 International licence, as recorded in the arXiv licence metadata for this exact version and shown on the abstract page https://arxiv.org/abs/2204.05442v7. Full rights notice: licenses/arxiv-2204.05442-CC-BY-4.0.txt.\par\smallskip

\noindent\textit{Editorial scope.} Counted unit: the Lemma whose source label is \texttt{lm:ApproxU} in the article (source0.tex lines 342--344, source label \texttt{lm:ApproxU}), delivered together with the article's own complete original proof of it (source0.tex lines 345--359, one \texttt{proof} environment of 15 lines). No mathematics is written, repaired or re-derived: statement and proof are the article's own text line for line. Required context retained uncounted: lm:Approximation (lines 334--336). Every definition, display and label of those lines is retained unchanged. Omitted from this package and not redistributed: 1--333, 337--341, 360--692. Nothing inside a retained excerpt is omitted or altered: every retained body is delivered exactly as the article's lines are written, so no omission receipt of any kind accompanies this package. Citations: the retained excerpts carry no citation command at all, so no bibliography entry of the article is transcribed and no reference is redistributed. Reference adaptation: none was needed and none was made; every reference inside the delivered unit resolves inside it, so no unresolved reference or citation is produced. Printed slips kept literal: no printed word, symbol, hypothesis or number of the counted statement or of its proof was altered. Layout and class: the article's own class is \texttt{amsart} with packages including amscd, amssymb, euscript, hyperref, fontenc; the wrapper is the lane's pinned \texttt{amsart} editorial wrapper at 10pt on A4, which restates the theorem-like environments the retained text needs and adds the article's own macro definitions that the retained text needs (\texttt{\textbackslash Frac}, \texttt{\textbackslash bU}, \texttt{\textbackslash cO}, \texttt{\textbackslash fp}), with extra wrapper packages (bm, bbm, dsfont, xspace, cancel, mathtools). The delivered numbering is the unit's own. Assets: no retained span contains a figure environment, a graphics-inclusion command, a PDF-page inclusion, a table or a TikZ picture; no image file is copied into this package, the package declares no asset dependency and no image file is redistributed with it. Licence: version arXiv:2204.05442v7 of this article is distributed under the Creative Commons Attribution 4.0 International licence, as recorded in the arXiv licence metadata for this exact version and shown on the abstract page https://arxiv.org/abs/2204.05442v7. A full rights notice is delivered at licenses/arxiv-2204.05442-CC-BY-4.0.txt. Characters: every retained excerpt is pure ASCII, so the delivered engine, XeLaTeX with its default Latin Modern text font, reports no missing character.
\begin{lemma}\label{lm:Approximation}
  Let $R$ be a Noetherian normal domain and $f\in R-\{0\}$. Let $\fp_1,\ldots,\fp_m\subset R$ be all the distinct minimal prime ideals containing $f$. Let $\overline\cO_j$ be the completion of the discrete valuation ring $R_{\fp_j}$ and let $\overline K_j$ be the fraction field of $\overline\cO_j$. Let $M$ be a finitely generated $R$-module. Assume that for all $j$ we are given $u_j\in M\otimes_R\overline K_j$. Then for all $N>0$ there are $h\in R$ and $u\in M\otimes_R R_{fh}$ such that for all $j$ we have $h\notin\fp_j$ and $u\equiv u_j\pmod{M\otimes_R\fp_j^N\overline\cO_j}$, where we view $u$ as an element of $M\otimes_R\overline K_j$ via the composition $R_{fh}\to\Frac(R)=\Frac(R_{\fp_j})\to\overline K_j$.
\end{lemma}

\begin{lemma}\label{lm:ApproxU}
  Let $R$, $f$, $\fp_j$, $\overline\cO_j$ and $\overline K_j$ be as above, let $\bU$ be an affine $R$-group scheme with a filtration $\bU=\bU_n\supset\bU_{n-1}\supset\ldots\supset\bU_0=\{e\}$ such that the successive quotients $\bU_i/\bU_{i-1}$ are isomorphic to additive group schemes of locally free $R$-modules of finite rank. Let $u_j\in\bU(\overline K_j)$ for $j=1,\ldots,m$. Then for $N>0$ there are $h\in R$ and $u\in\bU(R_{fh})$ such that for all $j$ we have $h\notin\fp_j$ and $u_j\in \bU^{(N)}_ju$.
\end{lemma}

\begin{proof}
  Let $l$ be an integer such that $0\le l\le n$. Assume that for all $j$ we have $u_j\in\bU_l(\overline K_j)$. We will show that for all $N>0$ we can find $h\in R$ and $u\in\bU_l(R_{fh})$ such that for all $j$ we have $h\notin\fp_j$ and $u_j\in (\bU_l)^{(N)}_ju$. When $l=n$, this statement is the claim of our lemma.

  We induct on $l$, the case $l=0$ being obvious. Assume that the statement is known for $l-1$. Let $\bar u_j\in(\bU_l/\bU_{l-1})(\overline K_j)$ be the image of $u_j$ in the quotient. By Lemma~\ref{lm:Approximation} we can find $h'\in R$ and $\bar w\in(\bU_l/\bU_{l-1})(R_{fh'})$ such that for all $j$ we have $h'\notin\fp_j$ and $\bar u_j=\bar v_j \bar w$, where $\bar v_j\in (\bU_l/\bU_{l-1})^{(N)}_j$.

  Since $H^1(R_{fh'},\bU_{l-1})=0$, we can lift $\bar w$ to an element $w\in\bU_l(R_{fh'})$. We claim that we can lift $\bar v_j$ to an element $v_j\in(\bU_l)^{(N)}_j$. Indeed, since $H^1(\overline\cO_j,\bU_{l-1})=0$, we can lift $\bar v_j$ to an element $v'\in\bU_l(\overline\cO_j)$. The image of $v'$ in $\bU_l(\overline\cO_j/\fp_j^N\overline\cO_j)$ is contained in $\bU_{l-1}(\overline\cO_j/\fp_j^N\overline\cO_j)$.
  Let $v''$ be a lift of this element to $\bU_{l-1}(\overline\cO_j)$. We can take $v_j:=v'(v'')^{-1}$.

  Then for all $j$ we have $u_j=v_jwu'_j$, where $u'_j\in\bU_{l-1}(\overline K_j)$. Since conjugation by any element is a continuous automorphism of $\bU_l(\overline K_j)$, we can find $M>0$ such that for all $j$ and all $a\in(\bU_l)^{(M)}_j$ we have $waw^{-1}\in(\bU_l)^{(N)}_j$. Applying the induction hypothesis, we find $h''\in R$ and $u'\in\bU_{l-1}(R_{fh''})$ such that for all $j$ we have $h''\notin\fp_j$ and $u'_j=a_ju'$, where $a_j\in(\bU_{l-1})^{(M)}_j$.

  Finally, take $h=h'h''$, we have
  \[
    u_j=v_jwu'_j=v_jwa_ju'=v_j(wa_jw^{-1})wu'\in(\bU_l)^{(N)}_j(wu').
  \qedhere\]
\end{proof}

\end{document}
